\chapter{Позначення}
\label{sec:notation}

Типи нейронів позначено першими чотирма літерами англійської назви головного медіатора: \nt{GLUT} --- глутамат, \nt{GABA} --- гамма-аміномасляна кислота, \nt{GLYC} --- гліцин, \nt{ACET} --- ацетилхолін, \nt{DOPA} --- дофамін, \nt{SERO} --- серотонін, \nt{HIST} --- гістамін, \nt{NORA} --- норадреналін, \nt{ADRE} --- адреналін, \nt{ASPA} --- аспартат (таблиця~\ref{tab:types}).

Літер менше, ніж величин, і в різних розділах одна літера іноді означає різне. У таблиці нижче кожне значення має свій рядок; у правому стовпці --- формула, де його введено.

\begin{center}
\footnotesize
\setlength{\tabcolsep}{4pt}
\renewcommand{\arraystretch}{1.12}
\begin{longtable}{>{\raggedright}p{2.0cm}>{\raggedright}p{9.6cm}>{\raggedright\arraybackslash}p{2.4cm}}
\toprule
Символ & Значення & Де введено \\
\midrule
\endfirsthead
\toprule
Символ & Значення & Де введено \\
\midrule
\endhead
\bottomrule
\endfoot
$a$ & крутість лінійної передавальної характеристики, $f(u)=au$ & \eqref{eq:feedback} \\
$a$ & рівень \nt{ACET}: $0$ --- читання, $1$ --- запис & \eqref{eq:acetnet} \\
$a_i$, $a$ & утома групи: у генераторі ритму; у гойдалці повільного сну & \eqref{eq:halfcentre}, \eqref{eq:updown} \\
$a_S$, $a_P$ & чутливість ритму серця до газу і до гальма & \eqref{eq:gasbrake} \\
$A(r)$, $A_0$ & відповідь синапсу на імпульс за частоти $r$; у відпочилого синапсу & \eqref{eq:depression} \\
$A_1$ & відповідь отримувача на одну порцію медіатора & \eqref{eq:quantal} \\
$A$ & площа мембрани нейрона разом із дендритами & \eqref{eq:dendriteload} \\
$A_f$, $A_s$ & частка поправки, що зберігається до наступного руху, у швидкого і повільного процесів & \eqref{eq:tworate} \\
$b_i$ & власний приплив ритмоводія & \eqref{eq:seroneuron} \\
$b$ & як швидко робота втомлює групу в генераторі ритму або в гойдалці сну & \eqref{eq:halfcentre}, \eqref{eq:updown} \\
$b$ & кількість біт на один відлік оцифрування & \eqref{eq:probedata} \\
$B_f$, $B_s$ & сила, з якою поправка вчиться на помилці, у швидкого і повільного процесів & \eqref{eq:tworate} \\
$c$ & концентрація медіатора в об'ємі & \eqref{eq:seroneuron} \\
$c$ & коефіцієнт кореляції шумів сусідніх нейронів & \eqref{eq:correlated} \\
$c$ & коефіцієнт, з яким різниця лічильників рухає ракетку & \eqref{eq:dishreadout} \\
$c_f$ & відповідь $C$-нейрона: найбільша з $s_{f,p}$ за всіма положеннями & \eqref{eq:scstage} \\
$c_m$ & питома ємність мембрани & \eqref{eq:dendriteload} \\
$C$ & ціна зусилля: дія за час $\tau_a$ коштує $C/\tau_a$ & \eqref{eq:vigor} \\
$d'$ & розрізненність двох входів: різниця, поділена на шум & \eqref{eq:gainsnr} \\
$d$, $d_j$ & затримка на шляху сигналу & \eqref{eq:glutneuron}, \eqref{eq:vstar} \\
$d$ & затримка петлі зворотного зв'язку з м'язом & \eqref{eq:delaystable} \\
$D$ & коефіцієнт дифузії & \eqref{eq:diffusionlength} \\
$D$ & знаменник у частотах \nt{GLUT}--\nt{GABA}-мережі & \eqref{eq:isn} \\
$D$ & час очікування відкладеної винагороди & \eqref{eq:patience} \\
$e$, $e_{ij}$ & слід у синапсі & \eqref{eq:threefactor} \\
$e$, $e_i$ & помилка виходу або руху, що надходить проводом помилки & \eqref{eq:lms}, \eqref{eq:adaptivefilter}, \eqref{eq:tworate} \\
$E_{\text{імп}}$ & енергія одного імпульсу разом із роботою синапсів & \eqref{eq:energy} \\
$E$ & частота \nt{GLUT}-групи & \eqref{eq:ei} \\
$E_E$ & рівноважний потенціал збуджувального синапсу & \eqref{eq:shunt} \\
$f$, $F$ & передавальна характеристика: частота як функція входу & \eqref{eq:ratenet}, \eqref{eq:gain} \\
$f$, $f_0$ & частота скорочень серця; власний ритм ритмоводія & \eqref{eq:gasbrake} \\
$f_s$ & частота оцифрування одного електрода & \eqref{eq:probedata} \\
$F(t)$, $F_1$ & сила м'яза; сила одного посмикування & \eqref{eq:forcesum} \\
$F_1$, $F_2$ & сили двох м'язів, що тягнуть суглоб у різні боки & \eqref{eq:antagonists} \\
$g$ & чутливість рецептора до концентрації & \eqref{eq:seroneuron} \\
$g_L$, $g_E$, $g_I$ & провідності витоку, збудження і гальмування & \eqref{eq:shunt} \\
$g$ & коефіцієнт підсилення, що його задає \nt{NORA} & \eqref{eq:gain}, \eqref{eq:machine} \\
$g$ & загальне гальмування від \nt{GABA} у мережі пам'яті & \eqref{eq:acetnet} \\
$g$ & підсилення воріт болю, що задається згори & \eqref{eq:gate} \\
$g$ & підсилення одного ступеня гормонального каскаду & \eqref{eq:cascade} \\
$g_j$ & фільтр $j$ зі своєю сталою часу на вході адаптивного фільтра & \eqref{eq:adaptivefilter} \\
$G$ & підсилення петлі зворотного зв'язку & \eqref{eq:feedback} \\
$G$ & швидкість виправлення в петлі із затримкою & \eqref{eq:delaystable} \\
$h$, $h_I$ & зовнішній вхід \nt{GLUT}-групи і \nt{GABA}-групи & \eqref{eq:ei}, \eqref{eq:isn} \\
$h(t)$ & одиночне посмикування м'яза & \eqref{eq:twitch} \\
$H(t)$ & верхній поріг тиску сну: дійшовши до нього, машина засинає & \eqref{eq:sleeppressure} \\
$I$, $I_i$ & зовнішній приплив до нейрона або групи & \eqref{eq:ratenet}, \eqref{eq:wta}, \eqref{eq:updown} \\
$I$ & частота \nt{GABA}-групи & \eqref{eq:ei} \\
$I$ & яскравість або інтенсивність подразника & \eqref{eq:photonnoise}, \eqref{eq:nakarushton} \\
$I$ & вхідний струм, що заряджає мембрану & \eqref{eq:dendriteload} \\
$k$ & кількість імпульсів у пачці & \eqref{eq:burst} \\
$k$ & у скільки разів приріст має перевищувати шум & \eqref{eq:photonnoise} \\
$k$ & коефіцієнт зворотного зв'язку гормонального каскаду & \eqref{eq:cascade} \\
$k_s$ & сила, з якою біль накачує сенсибілізацію & \eqref{eq:sensitization} \\
$K$ & кількість ознак результату & \eqref{eq:vectortd} \\
$K$ & жорсткість суглоба & \eqref{eq:antagonists} \\
$K$ & підсилення петлі гормонального каскаду, $K=g^3k$ & \eqref{eq:cascade}, \eqref{eq:cascadestable} \\
$K_\varphi$ & сила підлаштування добового генератора до зовнішнього сигналу & \eqref{eq:pll} \\
$L$ & довжина лінії затримок & \eqref{eq:itd} \\
$L$ & довжина проводу від датчика болю & \eqref{eq:paintime} \\
$L(t)$ & нижній поріг тиску сну: дійшовши до нього, машина прокидається & \eqref{eq:sleeppressure} \\
$M$ & міра причинного зв'язку провісника і винагороди & \eqref{eq:retro} \\
$M_k$ & очікування ознаки $k$ & \eqref{eq:vectortd} \\
$M_{ik}$ & сила гальмування від \nt{GABA}-нейрона $k$ & \eqref{eq:machine} \\
$M$ & обертальний момент суглоба & \eqref{eq:antagonists} \\
$M$ & кількість вхідних ліній змішувачів & \eqref{eq:cover} \\
$n$, $\bar n$ & кількість імпульсів за вікно; її середнє & \eqref{eq:rateerror} \\
$n$ & кількість входів порогового елемента & \eqref{eq:cover} \\
$n_\uparrow$, $n_\downarrow$ & лічба імпульсів ділянок «вгору» і «вниз» за вікно & \eqref{eq:dishreadout} \\
$N$ & кількість нейронів & \eqref{eq:correlated}, \eqref{eq:energy} \\
$N$ & кількість електродів щупа & \eqref{eq:probedata} \\
$N_s$ & кількість місць викиду між двома нейронами & \eqref{eq:quantal} \\
$p$ & імовірність викиду медіатора на імпульс & \eqref{eq:burst} \\
$p$ & частка активних нейронів у мережі пам'яті & \eqref{eq:acetnet} \\
$p$ & збурення, яке вчаться компенсувати & \eqref{eq:tworate} \\
$P$ & потужність, що витрачається на сигнали & \eqref{eq:energy}, \eqref{eq:dishpower} \\
$P$ & невизначеність збереженої оцінки & \eqref{eq:kalman} \\
$P$ & активність «гальма»: \nt{ACET}-проводів до серця & \eqref{eq:gasbrake} \\
$P_{\max}$ & скільки випадкових образів пороговий елемент може поділити на два класи & \eqref{eq:cover} \\
$P_{\text{пр}}$ & імовірність промаху & \eqref{eq:dishdensity} \\
$q$ & швидкість, з якою змінюється світ & \eqref{eq:kalman} \\
$q_k$ & частота \nt{GABA}-нейрона $k$ & \eqref{eq:machine} \\
$q_0$ & фонова активність гальмівного посередника & \eqref{eq:gate} \\
$q$ & частота гальмівного посередника у воротах болю & \eqref{eq:gate} \\
$q$ & у скільки разів наступна рухова одиниця сильніша за попередню & \eqref{eq:sizeprinciple} \\
$r$, $r_i$ & частота імпульсів нейрона & \eqref{eq:ratenet} \\
$r$, $r_t$ & винагорода (потік винагороди) & \eqref{eq:rpe}, \eqref{eq:ctd} \\
$\mathbf r(t)$ & вектор відгуків резервуара & \eqref{eq:reservoir} \\
$R$, $\bar R$ & отримана винагорода; її очікування або середнє за одиницю часу & \eqref{eq:perturbation}, \eqref{eq:vigor} \\
$R$ & дисперсія шуму на вході фільтра & \eqref{eq:kalman} \\
$R$, $r_0$ & відкладена і негайна винагороди & \eqref{eq:patience} \\
$R_{\max}$ & гранична швидкість передавання, біт/с & \eqref{eq:spikeentropy} \\
$r_{\max}$ & найбільша частота імпульсів & \eqref{eq:gain}, \eqref{eq:nakarushton} \\
$R_{\text{сирий}}$, $R_{\text{імп}}$ & потік даних щупа: сирий і лише імпульси, біт/с & \eqref{eq:probedata} \\
$s_j$ & знак виходів нейрона $j$ & \eqref{eq:dalesign} \\
$s_i$ & знак рецептора отримувача & \eqref{eq:seroneuron} \\
$s$, $\hat s$ & стан світу; його оцінка & \eqref{eq:rpe}, \eqref{eq:kalman} \\
$s_{f,p}$ & відповідь $S$-нейрона на ознаку $f$ у положенні $p$ & \eqref{eq:scstage} \\
$S$ & активність «газу»: \nt{NORA}-проводів до серця & \eqref{eq:gasbrake} \\
$S$ & тиск сну & \eqref{eq:sleeppressure} \\
$t_1$ & момент першого імпульсу & \eqref{eq:latency} \\
$t_k$ & момент $k$-го імпульсу & \eqref{eq:forcesum} \\
$t_{f,i}$ & зразок ознаки $f$ & \eqref{eq:scstage} \\
$T$ & вікно лічби імпульсів; час накопичення фотонів & \eqref{eq:rateerror}, \eqref{eq:photonnoise} \\
$T_c$ & час наростання посмикування м'яза & \eqref{eq:twitch} \\
$T_{\text{залп}}$ & інтервал між залпами культури & \eqref{eq:burstinterval} \\
$u$ & сумарний вхід нейрона & \eqref{eq:gain} \\
$u$, $u(t)$ & команда мозку: на вході адаптивного фільтра; гормональному каскаду & \eqref{eq:adaptivefilter}, \eqref{eq:cascade} \\
$u(t)$ & вхід резервуара & \eqref{eq:reservoir} \\
$U$ & рівень сталого входу & \eqref{eq:latency} \\
$U$ & частка запасу медіатора, що викидається на імпульс & \eqref{eq:depression} \\
$v$ & швидкість сигналу в проводі; швидкість руху плями & \eqref{eq:itd}, \eqref{eq:vstar} \\
$V$ & напруга на мембрані & \eqref{eq:glutneuron}, \eqref{eq:shunt} \\
$V$ & очікування майбутньої винагороди & \eqref{eq:rpe}, \eqref{eq:ctd} \\
$w$, $w_{ij}$, $W_{ij}$ & вага зв'язку & \eqref{eq:glutneuron}, \eqref{eq:ratenet}, \eqref{eq:acetnet}, \eqref{eq:lms}, \eqref{eq:adaptivefilter} \\
$\mathbf w_k$ & ваги входів $C$-нейрона $k$ & \eqref{eq:traceinvariance} \\
$w_{q\mu}$, $w_{q\nu}$, $w_{z\mu}$ & ваги зв'язків у воротах болю & \eqref{eq:gate} \\
$W_{\text{вих}}$ & ваги лінійного зчитування резервуара & \eqref{eq:reservoir} \\
$x$, $y$ & логічні входи і вихід & \eqref{eq:dualrail}, \eqref{eq:veto} \\
$x$, $y$ & активність відправника й отримувача & \eqref{eq:hebb} \\
$x$ & місце детектора на лінії затримок & \eqref{eq:itd} \\
$x_i$ & вхід від органів чуття & \eqref{eq:acetnet} \\
$x_p$ & вхід у положенні $p$ & \eqref{eq:scstage} \\
$\mathbf x$ & виходи $S$-нейронів, вхід $C$-нейрона & \eqref{eq:traceinvariance} \\
$x_j$ & вхід $j$ адаптивного фільтра: команда після фільтра $g_j$ & \eqref{eq:lms}, \eqref{eq:adaptivefilter} \\
$x_f$, $x_s$ & поправки швидкого і повільного процесів & \eqref{eq:tworate} \\
$x_1$, $x_2$, $x_3$ & рівні на трьох ступенях гормонального каскаду; $x_3$ --- кінцевий гормон & \eqref{eq:cascade} \\
$x^*$ & частка відновленого запасу медіатора, за якої починається нова лавина & \eqref{eq:burstinterval} \\
$y$ & зашумлений вхід фільтра & \eqref{eq:kalman} \\
$y$, $\hat y$ & сигнал датчиків; його передбачення адаптивним фільтром & \eqref{eq:adaptivefilter} \\
$y_k$, $\bar y_k$ & активність $C$-нейрона $k$; її слід & \eqref{eq:traceinvariance} \\
$y(t)$ & вихід зчитування резервуара & \eqref{eq:reservoir} \\
$z$ & частота виходу воріт болю & \eqref{eq:gate} \\
\midrule
$\alpha_i^\pm$ & вага додатних і від'ємних помилок & \eqref{eq:asymmetric} \\
$\beta$ & сила взаємного гальмування & \eqref{eq:wta} \\
$\beta$ & сила гальмування виходу воріт болю & \eqref{eq:gate} \\
$\gamma$ & коефіцієнт дисконтування & \eqref{eq:rpe} \\
$\delta(\cdot)$ & дельта-функція: миттєвий стрибок & \eqref{eq:glutneuron} \\
$\delta$, $\delta_t$ & помилка передбачення винагороди & \eqref{eq:rpe}, \eqref{eq:ctd} \\
$\delta t$ & різниця часів приходу звуку у два вуха & \eqref{eq:itd} \\
$\Delta$, $\Delta_{\max}$ & інтервал між імпульсами; вікно збігу & \eqref{eq:window} \\
$\Delta t$ & точність, з якою отримувач розрізняє час & \eqref{eq:spikeentropy} \\
$\Delta F$ & приріст сили від чергової рухової одиниці & \eqref{eq:sizeprinciple} \\
$\Delta I$ & помітний приріст яскравості & \eqref{eq:photonnoise} \\
$\Delta p$ & зсув ракетки за вікно & \eqref{eq:dishreadout} \\
$\Delta u$ & різниця входів двох груп & \eqref{eq:gainsnr} \\
$\Delta V$ & на скільки треба підняти напругу мембрани & \eqref{eq:dendriteload} \\
$\Delta x$ & відстань між сусідніми датчиками & \eqref{eq:vstar} \\
$\varepsilon$ & відносна різниця входів & \eqref{eq:wtagain} \\
$\eta$ & швидкість навчання & \eqref{eq:plasticity}, \eqref{eq:lms}, \eqref{eq:traceinvariance} \\
$\eta$ & шум у мережі після підсилювача & \eqref{eq:softmax} \\
$\mu$ & вхід дотику до воріт болю & \eqref{eq:gate} \\
$\nu$ & вхід болю & \eqref{eq:gate} \\
$\theta$ & поріг спрацювання & \eqref{eq:window}, \eqref{eq:scstage} \\
$\kappa$ & кількість медіатора на один імпульс & \eqref{eq:seroneuron} \\
$\kappa_i$ & відношення ваг додатних і від'ємних помилок & \eqref{eq:expectile} \\
$\kappa_m$ & зв'язок сили м'яза з його жорсткістю & \eqref{eq:antagonists} \\
$\ell$ & відстань, на яку розпливається медіатор & \eqref{eq:diffusionlength} \\
$\ell$ & плече, на якому м'яз тягне суглоб & \eqref{eq:antagonists} \\
$\xi$ & випадкове тремтіння виходу нейрона & \eqref{eq:perturbation} \\
$\rho(\boldsymbol\psi)$ & частка часу, проведеного в налаштуванні $\boldsymbol\psi$ & \eqref{eq:dishdensity} \\
$\sigma$ & розкид, масштаб шуму & \eqref{eq:rateerror}, \eqref{eq:softmax} \\
$\sigma$ & яскравість, за якої відповідь дорівнює половині найбільшої & \eqref{eq:nakarushton} \\
$\sigma_{\text{до}}$, $\sigma_{\text{після}}$ & шум до підсилювача і після нього & \eqref{eq:gainsnr} \\
$\tau$ & стала часу мембрани & \eqref{eq:glutneuron} \\
$\tau$ & стала часу ступеня гормонального каскаду & \eqref{eq:cascade} \\
$\tau_{\text{еф}}$ & стала часу групи, що збуджує сама себе & \eqref{eq:perfectint} \\
$\tau_c$ & час життя медіатора в об'ємі & \eqref{eq:seroneuron} \\
$\tau_E$, $\tau_I$ & сталі часу \nt{GLUT}- і \nt{GABA}-групи & \eqref{eq:ei} \\
$\tau_e$ & час життя сліду & \eqref{eq:threefactor} \\
$\tau_{\text{сл}}$ & час життя сліду активності $C$-нейрона & \eqref{eq:traceinvariance} \\
$\tau_s$ & час повернення чутливості після пошкодження & \eqref{eq:sensitization} \\
$\tau_r$, $\tau_d$ & час накопичення тиску сну під час неспання і його скидання уві сні & \eqref{eq:sleeppressure} \\
$\tau_{\text{відн}}$ & час відновлення запасу медіатора & \eqref{eq:depression} \\
$\tau_\gamma$ & горизонт планування & \eqref{eq:ctd} \\
$\tau_v$ & швидкість, з якою уточнюється очікування & \eqref{eq:ramp} \\
$\tau_a$ & час утоми в генераторі ритму або в гойдалці сну & \eqref{eq:halfcentre}, \eqref{eq:updown} \\
$\tau_a^*$ & найкращий час виконання дії & \eqref{eq:vigor} \\
$\Phi$ & права частина правила навчання & \eqref{eq:plasticity} \\
$\Phi$ & потік фотонів & \eqref{eq:photonnoise} \\
$\varphi_k$ & ознака $k$ отриманого результату & \eqref{eq:vectortd} \\
$\varphi$ & різниця фаз добового генератора і зовнішнього сигналу & \eqref{eq:pll} \\
$\boldsymbol\psi$ & налаштування мережі в моделі DishBrain & \eqref{eq:dishdensity} \\
$\omega$, $\omega_0$ & частота зовнішнього сигналу (світла); власна частота добового генератора & \eqref{eq:pll} \\
\end{longtable}
\end{center}
